% Caption for cand_c21_three.png.
% Every number verified against traces_torque3/ (curated, n=30) and
% g5grid/plane3_runs/ (plane, n=10); calibration in calib/ENERGY_CALIBRATION.md.

\caption{\textbf{Schedule choice decides task success on the two widowx scenes
and does nothing on the actuator-saturated one.}
%
\textbf{(A) The curated ladder.} Six schedules per workload, each at 30 seeds of
24 episodes, ordered reference-first and then by release period; ticks give the
schedule name over its measured release period in ms. Boxes are quartiles over
seeds, dots individual seeds, diamonds the mean; panel titles carry the units.
Across the ladder eggplant falls 54.9\%\,$\rightarrow$\,1.7\% and spoon
47.8\%\,$\rightarrow$\,1.4\% while their energy rises 1.71$\times$ and
1.42$\times$. Close drawer moves from 40.0\% to 36.8\% and its energy not at all
(1.00$\times$): \emph{no schedule differs significantly from the zero-latency
reference on a paired test over shared seeds} (every $p > 0.06$), so the flat
panel is a well-powered null rather than an absent measurement.
%
\textbf{(B) The full plane.} All 44 CP-SAT operating points at 10 seeds each.
Left: success against energy, normalised to each workload's own cheapest
schedule---the absolute figures sit in disjoint bands (spoon $\sim$90\,J,
eggplant $\sim$150\,J, close drawer $\sim$2.4\,kJ), so on a shared absolute axis
each series collapses to a vertical line and the within-workload ladder, which is
what is being compared, becomes unreadable. Right: success on the
(release period, deadline window) grid, one panel per workload on its own
gradient, with a rollout pair above it---success left, failure right. Stars mark
each workload's best schedule by success rate. Hatched cells are propagated from
a twin whose CP-SAT schedule is identical to within 0.04\,ms of latency and
0\,ms of cadence, so the simulator receives the same command; $\times$ and ?
mark cells that were infeasible and cells the solver did not resolve inside its
budget.
%
\textbf{Energy.} The raw $\int\!\sum\tau^2\mathrm{d}t$ channel is a real PD
drive-torque integral but cannot be converted to joules at face value: the
simulated torque runs about 76$\times$ what the same trajectory physically
requires (148\,N\,m rms against 1.94\,N\,m of gravity and Coriolis, saturating
20\% of substeps under a 200\,N\,m limit on a shoulder whose two XM430-W350
servos stall at 8.2\,N\,m combined). Plotted instead is the quasi-static lower
bound, $E = c\int\!\sum\tau_{\mathrm{grav}}^2\mathrm{d}t + P_{\mathrm{idle}}T$,
applied per episode and reduced as elsewhere (median over an arm's 24 episodes,
mean over seeds). For widowx, $c = 0.827\,\mathrm{W/(N\,m)^2}$ and
$P_{\mathrm{idle}} = 4.61$\,W follow from ROBOTIS' published stall torque and
current for the Dynamixel XM430-W350. \emph{Every constant here is a model rather
than a measurement}: the widowx values are derived from datasheets rather than
measured on the arm, and the Everyday Robots arm publishes no actuator data, so
its constant is set to give that arm the same dynamic/idle balance as the widowx
(0.810) at a literature-anchored 40\,W idle. Scaling a constant scales every
schedule within its workload equally, so the ratios in (A) and the normalised
axis in (B) are invariant to those choices; only absolute joules depend on them.
Derivation in \texttt{calib/ENERGY\_CALIBRATION.md}.}
